\section{Sharp Gaussian component moduli}
\label{sec:sharp-common-information}

The constants below refine the existing path certificate without using a
financial price as input. They are exact extrema for the stated class of
positive mass-one threshold measures. Their mathematical sharpness does
not establish priority in the literature. In particular, the two
component extrema need not be attained by the same perturbation; the
sum in the path certificate is not asserted to be jointly sharp.

Write $\phi(x)=(2\pi)^{-1/2}\exp(-x^2/2)$ and let $\Phi$ be its
distribution function. For $t\geq0$, define
\begin{equation}
 k(t)=\sup_{x\in\mathbb R}|\phi(x+t)-\phi(x)|,
 \qquad A(t)=2\Phi(t/2)-1.
 \label{eq:sharp-k-A}
\end{equation}
For $t,e\geq0$, put $r_t(x)=\Phi(x+t/2)-\Phi(x-t/2)$ and
\begin{equation}
 \Omega(t,e)=\sup_{x\in\mathbb R,\,|h|\leq e}
                    |r_t(x+h)-r_t(x)|.
 \label{eq:sharp-Omega-definition}
\end{equation}
Both $k(0)$ and $\Omega(t,0)=\Omega(0,e)$ are zero.

\begin{theorem}[Exact derivative coefficient]
\label{thm:sharp-common-derivative}
For $t>0$, there is exactly one $u_t>\max\{1,t/2\}$ satisfying
\begin{equation}
 u_t\tanh(tu_t/2)=t/2,
 \qquad
 k(t)=\phi(u_t-t/2)-\phi(u_t+t/2).
 \label{eq:sharp-k-root}
\end{equation}
It lies below $1+t/2$. The function $k$ is strictly increasing on
$(0,\infty)$ and
\begin{equation}
 0<k(t)<\min\{\phi(0),t\phi(1)\},\qquad
 \lim_{t\downarrow0}\frac{k(t)}{t}=\phi(1),\qquad
 \lim_{t\to\infty}k(t)=\phi(0).
 \label{eq:sharp-k-properties}
\end{equation}
For every positive Borel measure $\mu$ of mass one and every $\lambda>0$,
the smallest uniform coefficient, simultaneously over $m$ and $\mu$, in
$|\partial_m R_{\lambda,\mu}(m,d)|\leq\kappa(d,\lambda)$ is
\begin{equation}
 \kappa(d,\lambda)=\lambda^{-1}k(d/\lambda).
 \label{eq:sharp-common-coefficient}
\end{equation}
An atomic measure attains this coefficient.
\end{theorem}
\begin{proof}
For $a=t/2$ and $u>0$, Gaussian symmetry gives the positive function
$q_t(u)=\phi(u-a)-\phi(u+a)$; the absolute density difference reduces
to this function by translation and reflection. Its logarithmic
derivative is
\begin{equation}
 \frac{q_t'(u)}{q_t(u)}=-u+a\coth(au).
 \label{eq:sharp-k-log-derivative}
\end{equation}
The function $u\mapsto u\tanh(au)$ is continuous, strictly increasing
from zero to infinity. At $u=a$ its value is strictly below $a$, and at
$u=1$ its value is strictly below $a$, since $\tanh a<a$. Thus its
unique level-$a$ point exceeds both. The inequality
$\tanh x\geq x/(1+x)$ for $x\geq0$ follows from
$\exp(2x)\geq1+2x$. At $u=a+1$ it implies
\begin{equation}
 u\tanh(au)\geq
 \frac{a(a+1)^2}{1+a(a+1)}>a.
 \label{eq:sharp-k-upper-bracket}
\end{equation}
Hence $q_t$ increases up to $u_t$ and decreases thereafter, proving
the formula and bracket. At the maximiser $u_t>a$, differentiation
with respect to $t$ gives
\begin{equation}
 \partial_t q_t(u_t)=\tfrac12\{(u_t-a)\phi(u_t-a)
                              +(u_t+a)\phi(u_t+a)\}>0.
 \label{eq:sharp-k-monotonicity}
\end{equation}
The implicit function theorem applies because
$\partial_u[u\tanh(au)]>0$, so this is the derivative of $k$.
Both Gaussian densities are positive, giving $k(t)<\phi(0)$.
The fundamental theorem of calculus gives
$q_t(u)\leq t\sup|\phi'|=t\phi(1)$, strictly for a positive interval,
since the continuous function $|\phi'|$ attains its maximum only at
two isolated points. The bracket and root equation imply $u_t\to1$
as $t\downarrow0$, establishing the small-$t$ limit. Evaluating the
density difference at $x=0$ gives the lower bound
$\phi(0)-\phi(t)$, which proves the large-$t$ limit. Finally,
\begin{equation}
 \partial_m R_{\lambda,\mu}(m,d)=\lambda^{-1}
   \int\{\phi((m+d-a)/\lambda)-\phi((m-a)/\lambda)\}\,\mu(da).
 \label{eq:sharp-common-derivative-integral}
\end{equation}
Positivity and total mass one prove the coefficient bound. Choosing
$\mu=\delta_a$ and $m-a=\lambda(u_t-t/2)$ attains its absolute value.
\end{proof}

\begin{theorem}[Exact finite common-shift modulus]
\label{thm:sharp-finite-common-shift}
For $t,e>0$, put $b=(t+e)/2$ and $c=|t-e|/2$. There is exactly one
$u_{t,e}>\max\{t,e\}/2$ satisfying
\begin{equation}
 \log\cosh(bu_{t,e})-\log\cosh(cu_{t,e})=te/2.
 \label{eq:sharp-Omega-root}
\end{equation}
It is smaller than $\max\{t,e\}/2+1/\min\{t,e\}$, and
\begin{equation}
 \Omega(t,e)=r_t(u_{t,e}-e/2)-r_t(u_{t,e}+e/2).
 \label{eq:sharp-Omega-evaluation}
\end{equation}
The function $\Omega$ is symmetric in its arguments and strictly
increasing in each positive argument. Moreover,
\begin{equation}
 0<\Omega(t,e)<\min\{k(t)e,k(e)t,A(t),A(e)\}.
 \label{eq:sharp-Omega-comparison}
\end{equation}
For every positive mass-one $\mu$,
\begin{equation}
 |R_{\lambda,\mu}(m+h,d)-R_{\lambda,\mu}(m,d)|
 \leq\Omega(d/\lambda,E/\lambda),\qquad |h|\leq E, d\geq0.
 \label{eq:sharp-Omega-mixture}
\end{equation}
The right-hand side is the smallest uniform bound over this class,
and is attained by a point mass and a shift of magnitude $E$.
\end{theorem}
\begin{proof}
The function $r_t$ is even, strictly decreasing on $(0,\infty)$ and
positive everywhere for $t>0$. For a shift of magnitude $e$, reflection
reduces the absolute difference to
$D_{t,e}(u)=r_t(u-e/2)-r_t(u+e/2)$ for $u\geq0$. This quantity is
positive for $u>0$. Expanding the four distribution functions shows
$D_{t,e}=D_{e,t}$. Differentiation gives
\begin{equation}
 D_{t,e}'(u)=2\phi(u)\{
 e^{-c^2/2}\cosh(cu)-e^{-b^2/2}\cosh(bu)\}.
 \label{eq:sharp-Omega-derivative}
\end{equation}
Since $b>c\geq0$, the function
$H(u)=\log\cosh(bu)-\log\cosh(cu)$ is strictly increasing on
$(0,\infty)$: its derivative is
$b\tanh(bu)-c\tanh(cu)>0$, because $v\mapsto v\tanh(vu)$ is
strictly increasing for $v>0$. Also $H(0)=0$ and $H(u)\to\infty$.
Since $(b^2-c^2)/2=te/2$, the unique root in
\eqref{eq:sharp-Omega-root} gives the unique maximum of $D_{t,e}$.
At $u=e/2$, the derivative is $r_t'(0)-r_t'(e)>0$; symmetry also
gives positivity at $u=t/2$. This proves the lower bracket.
The inequalities $\cosh(bu)\geq e^{bu}/2$ and
$\cosh(cu)\leq e^{cu}$ imply
$H(u)\geq(b-c)u-\log2$. At the stated upper bracket this exceeds
$te/2$, since $b-c=\min\{t,e\}$ and $1>\log2$.

The maximum depends smoothly on $e>0$ by the implicit function
theorem. At its maximiser, $D_{t,e}'(u)=0$, so
$r_t'(u-e/2)=r_t'(u+e/2)$. Consequently
\begin{equation}
 \partial_e\Omega(t,e)
   =-r_t'(u_{t,e}-e/2)>0.
 \label{eq:sharp-Omega-monotonicity}
\end{equation}
The strict inequality uses $u_{t,e}>e/2$. It follows that the largest
modulus over all shifts of magnitude at most $e$ is attained at $e$.
Symmetry gives monotonicity in $t$ as well. Integrating $r_t'$ over
the interval of length $e$ gives $\Omega(t,e)<k(t)e$, strictly because
the continuous derivative cannot be constant at its extremal value on
a positive interval. Interchanging $t,e$ gives $\Omega(t,e)<k(e)t$.
The largest value of $r_t$ is $r_t(0)=A(t)$, whereas the smaller
of the two values in the attained difference remains positive; hence
$\Omega(t,e)<A(t)$ and symmetry gives $\Omega(t,e)<A(e)$.
For a fixed threshold $a$, the correction is
$r_{d/\lambda}((m-a+d/2)/\lambda)$. Integrating its common-shift
bound against $\mu$ proves \eqref{eq:sharp-Omega-mixture}.
An atomic measure with the lower centred argument
$u_{t,e}-e/2$ and shift $E$ attains the absolute difference.
\end{proof}

\begin{corollary}[Finite path certificate]
\label{cor:sharp-finite-path}
Under the premises of Theorem~\ref{thm:common-information-stability},
\begin{equation}
 |R_{\lambda,\mu}(\widehat m,\widehat\Delta)
                        -R_{\lambda,\mu}(m,\Delta)|
 \leq\Omega(\Delta/\lambda,E_m/\lambda)
                       +A(e_\Delta/\lambda).
 \label{eq:sharp-finite-path}
\end{equation}
The bound remains valid when the represented gap is negative and after
projection that fixes the true correction. A known upper true gap may
replace $\Delta$. If both gaps are nonnegative, the first argument of
$\Omega$ may instead use $\min\{\Delta,\widehat\Delta\}$.
\end{corollary}
\begin{proof}
Use the same two-segment path as the earlier theorem. The common-shift
segment receives \eqref{eq:sharp-Omega-mixture}. For the gap segment,
the integrand is the Gaussian probability of an interval of length
$|\widehat\Delta-\Delta|/\lambda$. The probability of an interval
of length $v$ is maximised by the centred interval, as follows by
differentiating $\Phi(x+v)-\Phi(x)$ and using Gaussian symmetry.
Its maximum is $A(v)$. This argument also covers a negative represented
gap. Positivity and mass one preserve the bound under integration.
The monotonicity of $A$ and $\Omega$ gives the stated replacements;
reversing the path gives the smaller represented-gap alternative.
Projection is $1$-Lipschitz and fixes the target.
\end{proof}

\subsection{Certified evaluation without a financial price input}

All inputs to $k$, $A$ and $\Omega$ are the smoothing scale and
independently certified error/gap bounds. The root equations have unique
roots, explicit brackets admitting rational outward enclosures and
strictly monotone left-hand sides. Directed interval evaluation and rational bisection
therefore enclose the roots and the constants to arbitrary precision.
When an interval evaluation cannot decide the sign, increasing arithmetic
precision decides it unless the rational trial point is the exact root;
that case is covered as follows. For a retained root bracket $[l,r]$,
a trial point $c\in[l,r]$, a certified residual $|F(c)|\leq\varepsilon$
and a certified derivative lower bound
$\inf_{[l,r]}F'\geq a>0$, the mean value theorem gives the new enclosure
\begin{equation}
 u_*\in[l,r]\cap[c-\varepsilon/a,c+\varepsilon/a].
 \label{eq:sharp-root-residual-enclosure}
\end{equation}
The analytic derivative is strictly positive on the stated lower brackets;
a sufficiently fine interval subdivision certifies a positive lower
bound on the compact retained bracket. Evaluating the expressions on the entire retained root interval
gives an upper bound, rather than treating a rounded root as exact.
Outward rational rounding may subsequently preserve the upper bound.
Neither a realised derivative price nor a simulated price enters this
procedure. The constants are applicable to every encoding satisfying the
same conditioning premise. They do not justify treating normal-label
dependent arithmetic perturbations as independent common shifts.

\subsection{The actual capped-payoff threshold measure}

For $f(x)=H^{-1}\min\{K(e^x-1)_+,H\}$, write
$\alpha=K/H$ and $\ell=\log(1+H/K)$. Its threshold measure is
\begin{equation}
 \mu_{\rm cap}(da)=\alpha e^a\mathbf1_{(0,\ell)}(a)\,da,
 \qquad\int\mu_{\rm cap}(da)=\alpha(e^\ell-1)=1.
 \label{eq:sharp-cap-measure}
\end{equation}
Its smoothed derivative, which is a density rather than a basket price,
has the explicit expression
\begin{equation}
 j_\lambda(m)=P_{\lambda,\mu_{\rm cap}}'(m)
 =\alpha e^{m+\lambda^2/2}
 \left\{\Phi\left(\frac{\ell-m-\lambda^2}{\lambda}\right)
              -\Phi\left(\frac{-m-\lambda^2}{\lambda}\right)\right\}.
 \label{eq:sharp-cap-density}
\end{equation}
Define the model-specific coefficient
\begin{equation}
 \kappa_{\rm cap}(d,\lambda)
  =\sup_{m\in\mathbb R}|j_\lambda(m+d)-j_\lambda(m)|.
 \label{eq:sharp-cap-kappa}
\end{equation}
For every $d>0$ and $\lambda>0$,
\begin{equation}
 0<\kappa_{\rm cap}(d,\lambda)
   <\kappa(d,\lambda),\qquad
 \kappa_{\rm cap}(d,\lambda)\leq(K+H)/H.
 \label{eq:sharp-cap-strict-improvement}
\end{equation}
To derive \eqref{eq:sharp-cap-density}, integrate
$\alpha e^a\lambda^{-1}\phi((m-a)/\lambda)$ over $(0,\ell)$
and complete the square in $a$. The result is
\eqref{eq:sharp-cap-density}. The function $j_\lambda$ is continuous,
strictly positive and tends to zero at both tails. Its translated
difference is not identically zero for $d>0$: otherwise
$j_\lambda(m+d)=j_\lambda(m)$ for every $m$, making this positive
function periodic, contrary to its zero tail limit. Thus its translated
difference attains its nonzero absolute maximum at a finite point.
For each fixed $m$, the Gaussian translated-density difference attains
its extremal absolute value at only two isolated threshold positions.
The positive absolutely continuous measure \eqref{eq:sharp-cap-measure}
gives those positions zero mass. Strictness of the integral triangle
bound at the attained maximiser proves the strict inequality in
\eqref{eq:sharp-cap-strict-improvement}. Finally,
$j_\lambda=\mathbb E f'(m+\lambda N)$ and
$0\leq f'\leq(K+H)/H$, so the difference of two such values has
absolute value at most $(K+H)/H$.

This coefficient is certifiably computable. For any rational $L>0$, all
$m\notin[-d-L\lambda,\ell+L\lambda]$ obey the tail bound
\begin{equation}
 |j_\lambda(m+d)-j_\lambda(m)|\leq\phi(L)/\lambda.
 \label{eq:sharp-cap-tail}
\end{equation}
Indeed, both Gaussian arguments then remain on the same side of zero
with magnitude at least $L$ for every threshold in $[0,\ell]$.
Their nonnegative densities each lie below $\phi(L)$, and their
absolute difference lies below the same bound. On the compact interval,
the integrand derivative gives
\begin{equation}
 |\partial_m\{j_\lambda(m+d)-j_\lambda(m)\}|
                         \leq2\phi(1)/\lambda^2.
 \label{eq:sharp-cap-mesh-Lipschitz}
\end{equation}
A mesh of maximum spacing $\delta$, with directed upper bounds $U_i$
for the absolute differences at its nodes and both endpoints included,
therefore yields the global certificate
\begin{equation}
 \kappa_{\rm cap}(d,\lambda)
 \leq\max\left\{\phi(L)/\lambda,
                       \max_i U_i+\phi(1)\delta/\lambda^2\right\}.
 \label{eq:sharp-cap-mesh-certificate}
\end{equation}
Letting $L$ increase, $\delta$ decrease and evaluation precision increase
converges to the exact supremum. A single coefficient at $d_*$ cannot be
used for every $0\leq d\leq d_*$ without a separate monotonicity argument.
A directly certified alternative is
\begin{equation}
 \overline\kappa_{\rm cap}(d_*,\lambda)
 =\sup_{0\leq v\leq d_*}\sup_m
                  |j_\lambda(m+v)-j_\lambda(m)|
 <\kappa(d_*,\lambda)\qquad(d_*>0).
 \label{eq:sharp-cap-uniform-gap-envelope}
\end{equation}
For its strict inequality, the joint function is continuous and tends
uniformly to zero in the two tails for $v\in[0,d_*]$. Its positive
maximum is attained on a compact set. At an attained pair with $v>0$,
the preceding strict cap inequality gives a value smaller than
$\kappa(v,\lambda)\leq\kappa(d_*,\lambda)$; at $v=0$ the value is zero.
To compute this envelope, retain the compact $m$ interval
$[-d_*-L\lambda,\ell+L\lambda]$ and a rectangular mesh with maximum
spacings $\delta_m,\delta_v$. Since the derivatives in the two
coordinates have absolute bounds $2\phi(1)/\lambda^2$ and
$\phi(1)/\lambda^2$, respectively, the maximum node upper bound plus
$\phi(1)(\delta_m+\delta_v/2)/\lambda^2$, combined with the tail
bound \eqref{eq:sharp-cap-tail}, is a global upper certificate.
Any use of this model-specific refinement
must preserve the represented conditioning information. In particular,
the shared-threshold implementation may have strike-label dependent
threshold errors. It cannot apply the integrated cap coefficient to
those errors. The strike-independent common endpoint shift may use the
cap coefficient; the remaining strike-dependent shift must retain a
uniform fixed-threshold allowance. Raw-payoff and numerical encodings
may use the cap coefficient only after their common endpoint errors have
been separated from scalar-map errors and inner-label dependent
remainders as in the existing guarded certificates.

\subsection{Finite remainders that depend on the inner label}

These allowances remain separate from common-information smoothing.
For any Borel set $B$ of Lebesgue measure at most $w\geq0$,
\begin{equation}
 \int_B\phi(x)\,dx\leq A(w).
 \label{eq:sharp-normal-set-concentration}
\end{equation}
Indeed, let $I=[-w/2,w/2]$. The density is at most $\phi(w/2)$ on
$B\setminus I$ and at least that value on $I\setminus B$. Thus
\begin{equation}
 \int_B\phi-\int_I\phi
 \leq\phi(w/2)\{|B\setminus I|-|I\setminus B|\}\leq0.
 \label{eq:sharp-normal-set-proof}
\end{equation}
The two boundary strips in a Gaussian interval expansion or shrinkage
by $E_s/\lambda$ have total length at most $2E_s/\lambda$.
If shrinkage empties the interval, its original length is at most this
same bound. The existing interval sandwich therefore sharpens the
label-dependent shared-threshold allowance to
\begin{equation}
 2D_N+A(2E_s/\lambda)
       \leq2D_N+2\phi(0)E_s/\lambda.
 \label{eq:sharp-threshold-inner-remainder}
\end{equation}
Only the continuous strip probability is bounded by concentration;
the finite law still uses its retained interval CDF discrepancy and
strict/left endpoint conventions. This argument imposes no independence
between the perturbation and the inner normal label.

The normalised cap has the exact finite payoff modulus
\begin{equation}
 \omega_f(h)=\sup_x|f(x+h)-f(x)|
       =\min\{1,(\alpha+1)(1-e^{-|h|})\}.
 \label{eq:sharp-cap-payoff-modulus}
\end{equation}
For $h\geq0$, the increment integrates $\alpha e^x$ over the intersection
of an interval of length $h$ with $(0,\ell)$. Since that density is
increasing inside its support, its largest integral uses the rightmost
subinterval of length $\min\{h,\ell\}$. This proves the formula;
negative shifts give the same modulus by translation. For $d\geq0$,
the difference $g_d(x)=f(x+d)-f(x)$ lies in $[0,v(d)]$, where
$v(d)=\min\{1,(\alpha+1)(1-e^{-d})\}$. If $h\geq0$, both increments
$f(x+d+h)-f(x+d)$ and $f(x+h)-f(x)$ lie in $[0,\omega_f(h)]$.
Their difference consequently has magnitude at most $\omega_f(h)$.
The same argument follows for $h<0$ by translation. Hence
\begin{equation}
 |g_d(x+\zeta)-g_d(x)|
 \leq\min\{v(d),\omega_f(E_s)\}\qquad(|\zeta|\leq E_s).
 \label{eq:sharp-raw-inner-remainder}
\end{equation}
This replaces the former $(\alpha+1)E_s$ allowance and is valid
pointwise, including arbitrary dependence on the inner label.
